\section{Numerical Sensitivity Study}
\label{sec:sensitivity}

Three scheme parameters were varied in controlled experiments to establish whether
the reported forces depend on them. All variants were built from a separate copy of
the source tree, so the production build was never modified; the baseline variant
binary is byte-identical to the unmodified build of the same tree, which confirms the
instrumentation compiles away at the baseline value and so cannot have shifted the
reference point of the sweep. It is not evidence about the instrumentation at other
values, where the code does differ by construction. Every variant was run one at a
time under the CPU quota of \cref{sec:quota}.

\paragraph{Which revisions these variants were built from.} The twelve sweep runs predate
the submitted binary and it is worth being specific rather than saying only ``a separate copy
of the source tree'', which a reader would reasonably take to mean the same revision. The
dissipation-floor group was built at \texttt{97f37254d00c}, the spatial-order and
limiter-parameter groups at \texttt{26ef1cb80b86}, and the submitted production results at
\texttt{cc42ab4facfd}.

What licenses the comparison is that the intervening commits did not touch the discretization
these sweeps probe. That argument is strongest for the floor group, where the baseline binary
was shown byte-identical to an unmodified build of its own tree, and weaker for the
order and limiter groups, where no such control was run. The conclusions those two groups
support are correspondingly coarse --- a $60\,\%$ drag difference between first and second
order, and a $0.81\,\%$ spread over a decade of the limiter constant --- and neither depends
on a digit that a revision difference could plausibly move. They are reported as
order-of-magnitude sensitivities, which is all they are used for.

\subsection{The linear-wave dissipation floor is unnecessary for these cases}
\label{sec:floorsens}

\Cref{sec:lineardissipation} disclosed that the Roe flux applies a dissipation
floor of $\pLinearFloor$ times the local maximum wave speed to \emph{all} waves,
including the entropy and shear waves whose eigenvalue $u_n$ vanishes on a
stagnation streamline. That is added artificial dissipation beyond the classical
Harten--Hyman entropy fix, and it was introduced prophylactically to suppress a
checkerboard mode. The controlled experiment does not support the necessity of that
choice, and the result is reported here as it came out.

\begin{table}[htbp]
  \centering
  \caption{Cylinder $\Reinf=20$, converged runs, with the linear-wave floor
  coefficient swept. All three variants converge to 6.73 residual orders.}
  \label{tab:floorsweep}
  \small
  \begin{tabular}{rrrr}
    \toprule
    Floor coefficient & $C_D$ & Pressure & Viscous \\
    \midrule
    \num{0.05} (production) & \num{2.018119} & \num{1.220943} & \num{0.797176} \\
    \num{0.01} & \num{2.017701} & \num{1.220480} & \num{0.797221} \\
    \num{0.00} (disabled) & \num{2.017754} & \num{1.220447} & \num{0.797307} \\
    \bottomrule
  \end{tabular}
\end{table}

\Cref{tab:floorsweep} shows a total spread of $\num{4.2e-4}$ in $C_D$, which is
$0.021\,\%$ of the coefficient --- far below any physically meaningful difference and
below the case's own convergence tolerance. Disabling the floor entirely produced no
checkerboard mode by three independent measures: a second-to-first surface
difference ratio of $0.128$, four sign alternations per 100 wall faces, and mirror
asymmetry of $\num{6e-4}$.

The cylinder is a mild test, so the experiment was repeated on the case designed to
provoke the mode: \texttt{naca0012\_m015\_inviscid}, where $u_n$ vanishes exactly on
the stagnation streamline of a symmetric body at zero incidence and there is
\emph{no physical viscosity} to damp a carbuncle. Over a fixed \num{8000}-step
budget, restricted to the leading-edge region $x/c<0.1$:

\begin{table}[htbp]
  \centering
  \caption{NACA0012 $\Minf=0.15$ inviscid stagnation-region diagnostics at a fixed
  \num{8000}-step budget, leading edge $x/c<0.1$. The exact stagnation value is
  $C_p=1$.}
  \label{tab:stagnation}
  \small
  \begin{tabular}{lrr}
    \toprule
    Diagnostic & Floor \num{0.05} & Floor \num{0.00} \\
    \midrule
    $\mathrm{d}C_p$ sign alternations & 1 of 119 & 1 of 119 \\
    Oscillation ratio $d_2/d_1$ & \num{0.263} & \num{0.274} \\
    Stagnation $C_p$ & \num{1.002635} & \num{1.002502} \\
    Mirror asymmetry & \num{3.42e-3} & \num{3.28e-3} \\
    \bottomrule
  \end{tabular}
\end{table}

The two columns of \cref{tab:stagnation} are indistinguishable, and
\cref{fig:stagnation} shows the surface distributions overlaying. Where the
convergence behaviour differs at all it mildly \emph{favours} the disabled floor:
171 versus 168 CFL rejections, a residual-increase fraction of $50.5\,\%$ versus
$50.7\,\%$, and maximum inner iterations of 33 versus 19.

\begin{figure}[htbp]
  \centering
  \cnsfig{naca_stagnation_cp.png}{\textwidth}
  \caption{NACA0012 $\Minf=0.15$ inviscid: surface $C_p$ and its differences with
  the linear-wave floor enabled (\num{0.05}) and disabled (\num{0}), from the
  \texttt{surface.csv} of each variant run. The two are visually
  indistinguishable, including in the stagnation region where a carbuncle would
  appear.}
  \label{fig:stagnation}
\end{figure}

\paragraph{Conclusion, stated against the design choice.} For the eight benchmark
cases the floor is \textbf{a conservative but unnecessary safeguard}. It changes no
reported number materially, and removing it produces no checkerboard mode even in
the case constructed to provoke one. It is retained as insurance for meshes and
conditions not exercised here, but \emph{no result in this report depends on it}.

\paragraph{Scope limit.} This conclusion does not generalise. Neither test
exercises the regime in which the Roe carbuncle is most notorious --- a strong bow
shock at high Mach number on shock-aligned quadrilateral cells. The $\Minf=0.80$
and $\Minf=2.0$ cases would be where to look, and that experiment was not
performed. What is established is that the floor is unnecessary \emph{for these
cases}, not that it is unnecessary in general.

\paragraph{A diagnostic that had to be localised.} One statistic from this study is
worth recording as a caution. Computed over the whole surface, the oscillation
ratio is $1.72$, which looks alarming. It is not oscillation: it comes entirely
from the blunt trailing edge at $x/c=1.005$, where the $C_p$ jump is about $1.46$.
It appears mirror-symmetrically on both surfaces, in both variants, and in the
converged production run ($1.728$), and it falls to $0.263$ when the final $1\,\%$
of chord is excluded. Every surface-oscillation statistic quoted above therefore
excludes the trailing edge, and is labelled as doing so. A diagnostic that is not
localised before being interpreted can easily manufacture a defect that does not
exist.

\subsection{Second-order reconstruction: what it is worth}
\label{sec:ordersens}

Running the cylinder $\Reinf=20$ case first order instead of second gives
$C_D=\num{3.239634}$ against the second-order $\num{2.018119}$: first order
\textbf{overpredicts drag by $60.5\,\%$}, and almost all of the error is in the
pressure component, which rises by roughly $90\,\%$. Since the literature value is
$2.0$--$2.1$, the first-order result is not merely different but wrong, while the
second-order result is in range.

This is the most direct evidence available that the piecewise-linear
reconstruction of \cref{sec:reconstruction} does substantial work rather than
being nominally present. It also quantifies what the numerical dissipation of a
first-order scheme costs on this mesh: excessive dissipation thickens the
separated wake, reduces base-pressure recovery, and inflates form drag.

\subsection{Limiter parameter}
\label{sec:limitersens}

Varying the Venkatakrishnan constant $K$ over a full decade gives
$C_D=\num{2.032244}$ at $K=1$, $\num{2.018119}$ at $K=\pVenkatK$ (production), and
$\num{2.015887}$ at $K=10$. The variation is monotone with a total spread of
$0.81\,\%$, so the reported drag is not sensitive to this parameter at the level
that matters for the conclusions here. The monotone trend is the expected
direction: a larger $K$ widens the smooth-region threshold of
\cref{eq:limiter}, limits less, and returns slightly less dissipative results.

\subsection{Independent reproduction}
\label{sec:reproduction}

The variant binaries were configured and built from a separate copy of the source
tree through an independent CMake configure. The floor-$\pLinearFloor$ variant
reproduces $C_D=\num{2.018119}$ for the cylinder $\Reinf=20$ case, agreeing to seven digits
with a production-configuration probe run under the same conditions --- second-order,
terminating at step $2783$. This is an independent reproduction of the cylinder result
rather than a re-reading of the same output files.

One clarification so the comparison can be checked: that seven-digit agreement is against
the matched probe, \emph{not} against the submitted run in \cref{tab:forces}. The submitted
run converges at step $\cnsStepsCylSteady$ to $C_D=\cnsCdCylSteady$, a slightly different
point on the same decaying tail, so a reader comparing $2.018119$ against the results table
will find agreement to four digits rather than seven. Both statements are true of different
run lengths, and \cref{sec:digits} explains why the fourth digit is where they part company:
the cylinder's tail still has $\cnsTailRemainingCylSteady$ of movement left at termination.

All force values quoted in this report are machine-extracted from the submitted
CSV and JSON outputs rather than transcribed by hand --- both for the tables
generated by \texttt{report/harvest\_numbers.py} and for the sensitivity values
above, which were verified against generated JSON. That verification caught one
transcription error in a draft of the sensitivity summary, which is the reason the
practice is worth following.
